\section{Related Work}
\label{sec:related_work}

This section reviews the work needed to position this thesis: image-only chest X-ray classification, multimodal clinical fusion, temporal/prior-aware modelling, class-imbalance methods, CXR-LT challenge solutions, and the CaMCheX framework.

\subsection{Single-Modal Chest X-Ray Classification}
\label{subsec:related_single_modal}

Public datasets such as the NIH ChestX-ray8/ChestX-ray14 dataset, CheXpert, and MIMIC-CXR have enabled large-scale chest X-ray classification research by providing radiographs with labels extracted from radiology reports \cite{wang2017chestxray8,irvin2019chexpert,johnson2019mimiccxr}. Convolutional neural networks such as ResNet and DenseNet have been commonly used because they learn hierarchical visual features from medical images while remaining practical to train \cite{he2016resnet,huang2017densenet}. DenseNet-based models are particularly prevalent because dense connections encourage feature reuse and improve gradient flow. More recently, modern visual backbones such as Vision Transformers, Swin Transformers, ConvNeXt, and ConvNeXt-V2 have been explored for medical image analysis. ConvNeXt-V2 combines convolutional design with modern architectural improvements and provides a strong visual feature extractor while remaining more efficient than very large transformer-based models \cite{woo2023convnextv2}. This motivates the use of a lightweight modern convolutional backbone as the visual encoder in this thesis.

However, image-only classifiers treat the radiograph as an isolated input. This differs from clinical practice, where radiologists interpret images together with indications, symptoms, vital signs, and prior studies. Image-only models are therefore important baselines, but they cannot use the full clinical context available during real chest X-ray interpretation.

\subsection{Multimodal Clinical Fusion}
\label{subsec:related_multimodal}

Multimodal learning aims to combine information from different sources in order to produce richer and more clinically meaningful representations. In radiology, the most common multimodal setting is image-text learning, where chest X-rays are paired with radiology reports. Since radiology reports describe visual findings and clinical impressions, they provide a natural source of supervision for learning image-text representations.

Prior work often uses contrastive learning to align medical images and text. ConVIRT learns medical visual representations by contrasting paired radiographs and reports, showing that naturally occurring image-report pairs can improve representation learning without requiring additional manual annotations \cite{zhang2022convirt}. GLoRIA extends this idea by learning both global and local image-text representations, aligning image regions with words or phrases from radiology reports \cite{huang2021gloria}. BioViL further improves biomedical vision-language pretraining by using stronger domain-specific text modelling through Chest X-Ray Bidirectional Encoder Representations from Transformers (CXR-BERT) and radiology-specific semantic learning \cite{boecking2022biovil}. MedCLIP relaxes the requirement for strictly paired image-text data by using unpaired medical images and text with a semantic matching objective \cite{wang2022medclip}.

These methods show the value of image-text learning, but many focus on pretraining, retrieval, phrase grounding, or report-based alignment. They do not always model the clinical context available before interpretation, such as indication text, vital signs, and prior studies. In addition, radiology reports are usually written after image interpretation, so using reports as direct input for current-study prediction can create an unrealistic setting if the goal is to support interpretation before the final report is produced.

CaMCheX is the closest prior work to this thesis because it moves beyond image-report pretraining and directly uses clinical context for chest X-ray classification \cite{sloan2025camchex}. Section~\ref{subsec:related_camchex_results} discusses CaMCheX separately because it is both a methodological reference and a benchmark reference for this work.

\subsection{Temporal and Prior-Aware Medical Modelling}
\label{subsec:related_temporal}

Temporal information is important in radiology because interpretation often depends on comparison with prior studies. A finding on a current chest X-ray may be clinically interpreted as new, stable, improved, or worsened depending on previous examinations. For example, a pleural effusion that is increasing over time may require different clinical attention from one that is unchanged from a prior study. Therefore, prior imaging and historical clinical information can provide useful context for current-study prediction.

Recent biomedical vision-language work has started to incorporate temporal structure. BioViL-T explicitly models prior images and reports when they are available, showing that temporal context can improve downstream tasks such as progression classification, phrase grounding, and report generation \cite{bannur2023biovilt}. The associated MS-CXR-T benchmark further highlights the importance of evaluating whether models understand temporal disease progression, including whether findings are improving, stable, or worsening \cite{bannur2023biovilt}.

Prior-aware modelling also introduces practical risks. Some patients have no prior study, and the time gap between studies can vary from days to years. Historical reports or labels can also create shortcut learning if they are used carelessly. For this reason, this thesis treats prior information as auxiliary context rather than as the main evidence, and uses time-gap embeddings and prior dropout to reduce over-dependence on historical information.

\subsection{Mitigating Class Imbalance in Medical Datasets}
\label{subsec:related_class_imbalance}

Chest X-ray classification is naturally a multi-label problem because a single study may contain multiple findings. It is also long-tailed because some abnormalities are common, while others are rare. The CXR-LT challenge highlights this issue by evaluating models on long-tailed multi-label chest X-ray classification, where rare findings and label co-occurrence make the task more difficult than standard balanced classification \cite{holste2024cxrlt,holste2025cxrlt2024}.

A common baseline loss for multi-label classification is binary cross-entropy (BCE), where each label is treated as an independent binary prediction. Although simple and widely used, BCE can be dominated by frequent negative labels. In chest X-ray datasets, most labels are negative for most studies, so the model may learn to focus on easy negatives rather than rare positive findings. This can reduce sensitivity for minority classes, which is problematic in medical settings where rare findings may still be clinically important.

Focal Loss addresses class imbalance by reducing the contribution of easy examples and focusing training on harder examples \cite{lin2017focalloss}. However, standard focal loss was originally designed for dense object detection and does not fully address the asymmetric nature of multi-label classification, where the number of negative labels is usually much larger than the number of positive labels.

Asymmetric Loss (ASL) was designed for multi-label classification and is more suitable for this setting \cite{ridnik2021asymmetric}. It applies different focusing behaviour to positive and negative samples, so easy negative labels are down-weighted more strongly while positive labels remain informative. This is useful for long-tailed chest X-ray classification, where rare positive findings can otherwise be overwhelmed by the large number of negative labels.

\subsection{CXR-LT 2023 Challenge Solutions}
\label{subsec:related_cxrlt_solutions}

Since our task follows long-tailed multi-label chest X-ray classification, the CXR-LT 2023 challenge provides useful reference methods and performance levels \cite{holste2024cxrlt}. Built on the MIMIC-CXR database, the challenge evaluated models on 26 findings—including 12 rare classes extracted from clinical reports—using patient-level splits to avoid leakage, with mean Average Precision (mAP) as the primary metric and mean AUROC (mAUROC) as an auxiliary metric \cite{holste2024cxrlt,holste2023cxrlt_physionet}. We focus on four representative submissions: the top-performing CheXFusion model (which employs a Multi-Label Decoder (ML-Decoder) head) \cite{kim2023chexfusion}, the second-ranked ensemble system \cite{nguyenmau2023advanced}, the robust asymmetric loss approach by Park et al. \cite{park2023robust}, and the text-aligned bag-of-tricks framework by Hong et al. \cite{hong2023bagoftricks}.


\begin{table}[H]
    \centering
    \caption[Representative CXR-LT 2023 challenge solutions]{Representative CXR-LT 2023 challenge solutions. Official aggregate test-phase results and rankings are summarized from the CXR-LT challenge overview \cite{holste2024cxrlt}, while individual method descriptions are cited in each row.}
    \label{tab:cxrlt2023_results}
    \begin{adjustbox}{width=\textwidth}
    \begin{tabular}{clcc}
        \hline
        Rank & Team / Strategy & mAP & mAUROC \\
        \hline
        1 & CheXFusion \cite{kim2023chexfusion} (ConvNeXt-S, multi-view Transformer, ML-Decoder, weighted ASL) & 0.372 & 0.850 \\
        2 & Nguyen-Mau et al. \cite{nguyenmau2023advanced} (EfficientNetV2/ConvNeXt, test-time augmentation (TTA)) & 0.354 & 0.836 \\
        4 & Robust Asymmetric Loss \cite{park2023robust} (ConvNeXt-B, Robust Asymmetric Loss) & 0.351 & 0.838 \\
        5 & Bag of Tricks \cite{hong2023bagoftricks} (ResNet50, text encoder, external data, TTA/ensemble) & 0.349 & 0.836 \\
        \hline
    \end{tabular}
    \end{adjustbox}
\end{table}

The challenge results also show a clear gap between common and rare findings. Table~\ref{tab:cxrlt2023_longtail} reports mAP by label-frequency group for the selected challenge teams. CheXFusion achieved 0.499 mAP on head classes but only 0.242 mAP on tail classes. Similar gaps appear across the other submissions, showing that rare-label recognition remains difficult even for strong challenge systems.

\begin{table}[H]
    \centering
    \caption[CXR-LT 2023 performance by label-frequency group]{CXR-LT 2023 performance by label-frequency group for representative teams. The grouped mAP values are summarized from the CXR-LT challenge overview \cite{holste2024cxrlt}.}
    \label{tab:cxrlt2023_longtail}
    \begin{tabular}{lcccc}
        \hline
        Team & Overall mAP & Head mAP & Medium mAP & Tail mAP \\
        \hline
        CheXFusion \cite{kim2023chexfusion} (Rank 1) & 0.372 & 0.499 & 0.246 & 0.242 \\
        Nguyen-Mau et al. \cite{nguyenmau2023advanced} (Rank 2) & 0.354 & 0.482 & 0.226 & 0.227 \\
        Park et al. \cite{park2023robust} (Rank 4) & 0.351 & 0.480 & 0.221 & 0.220 \\
        Hong et al. \cite{hong2023bagoftricks} (Rank 5) & 0.349 & 0.474 & 0.218 & 0.243 \\
        \hline
    \end{tabular}
\end{table}

These results motivate studying whether clinical context and prior-study information can help long-tail classification. Instead of focusing on leaderboard-style ensembling, this thesis evaluates a resource-constrained multimodal model that uses indication text, vital signs, and prior-study context.

\subsection{CaMCheX Results and Relevance to This Thesis}
\label{subsec:related_camchex_results}

CaMCheX is the primary reference for this work because it combines multi-view chest radiographs with clinical indications and vital signs using a multimodal transformer-based framework \cite{sloan2025camchex}. As summarized in Table~\ref{tab:camchex_reference_results}, CaMCheX reported 0.576 mAP and 0.916 mAUROC on CXR-LT 2023, and 0.725 mAP and 0.904 mAUROC on CXR-LT 2024 Task 2. These results suggest that clinical context can improve chest X-ray classification, although they should be interpreted as external benchmark results rather than controlled comparisons with this thesis. We include these benchmark results as context only, because our dataset scale, resolution, and training resources differ from the published CaMCheX setting.

\begin{table}[H]
    \centering
    \caption[Reference results reported by CaMCheX on classification benchmarks]{Reference results reported by CaMCheX on chest X-ray classification benchmarks. Published results are shown as external benchmark references using the same evaluation metrics. Differences in training scale, preprocessing, resolution, and pretraining should be considered when interpreting the comparison.}
    \label{tab:camchex_reference_results}
    \begin{tabular}{lccc}
        \hline
        Benchmark & Model & mAP & mAUROC \\
        \hline
        CXR-LT 2023 & CheXFusion & 0.372 & 0.850 \\
        CXR-LT 2023 & CaMCheX & 0.576 & 0.916 \\
        CXR-LT 2024 Task 2 & Best CXR-LT 2024 competitor & 0.526 & 0.833 \\
        CXR-LT 2024 Task 2 & CaMCheX & 0.725 & 0.904 \\
        MIMIC-CXR labels & CaMCheX & 0.669 & 0.934 \\
        \hline
    \end{tabular}
\end{table}

As shown in Table~\ref{tab:camchex_longtail_results}, the performance of CaMCheX varies across frequency groups, achieving high mAP on head classes but lower mAP on medium and tail classes. This pattern confirms that long-tailed classification remains difficult even for strong multimodal models.

\begin{table}[H]
    \centering
    \caption{CXR-LT 2023 frequency-group mAP for CheXFusion and CaMCheX.}
    \label{tab:camchex_longtail_results}
    \begin{tabular}{lcccc}
        \hline
        Model & Overall mAP & Head mAP & Medium mAP & Tail mAP \\
        \hline
        CheXFusion & 0.372 & 0.499 & 0.246 & 0.242 \\
        CaMCheX & 0.576 & 0.786 & 0.550 & 0.408 \\
        \hline
    \end{tabular}
\end{table}

These results motivate several design choices in this thesis, including multimodal fusion and asymmetric loss. However, this thesis does not attempt an exact reproduction of CaMCheX. The published CaMCheX setting uses different training scale, preprocessing, resolution, and pretraining conditions. This thesis instead implements a resource-constrained adaptation and adds prior-study temporal context.
